\section{Test Case 6: Discovery-Mode Analysis of Galaxy Properties}

The first five test cases in this paper demonstrate ASTRA's ability to correctly recover known astrophysical results using real observational data. This validates that ASTRA's integrated framework can perform established analysis tasks while maintaining proper uncertainty quantification and physical validation.

\textbf{This sixth test case addresses a different question}: can ASTRA discover new patterns and identify interesting objects in galaxy survey data without being told what to look for? We analyze a realistic galaxy sample based on published SDSS scaling relations, testing whether ASTRA can (1) recover known astrophysical relations, (2) identify outliers and unusual objects, and (3) generate testable hypotheses for follow-up observations.

\subsection{The Galaxy Discovery Problem}

Modern galaxy surveys measure hundreds of properties for millions of galaxies. While established scaling relations like the star-forming main sequence and mass-metallicity relation are well-known, these surveys also contain:
\begin{itemize}
\item \textbf{Starburst galaxies}: Galaxies with elevated star formation rates, potentially triggered by mergers or interactions
\item \textbf{Post-starburst galaxies}: Galaxies that recently quenched star formation, potentially due to major mergers or AGN feedback
\item \textbf{Transitioning galaxies}: Objects between star-forming and quenched states
\item \textbf{Outliers}: Galaxies that deviate significantly from established relations
\end{itemize}

Identifying these objects programmatically is challenging because they are rare (typically $<5$\% of the sample) and can only be distinguished by analyzing multiple properties simultaneously in multi-dimensional space.

\subsection{Data and Methods}

We analyzed a sample of 3,000 galaxies based on published SDSS scaling relations from Kauffmann et al. (2003), Tremonti et al. (2004), and Brinchmann et al. (2004). The sample includes galaxies in the redshift range $0.01 < z < 0.15$ with measurements of:
\begin{itemize}
\item Stellar mass ($\log M_*/M_\odot$)
\item Star formation rate (log SFR)
\item Gas-phase metallicity (12 + log O/H)
\item Optical color ($g-r$)
\item Structural parameters (Petrosian radius, axis ratio)
\item Specific star formation rate (sSFR = SFR/$M_*$)
\end{itemize}

The data is generated using empirical relations from the literature to produce realistic SDSS-like galaxy populations, including a small fraction ($\sim5$\%) of starburst, post-starburst, and transitioning galaxies for discovery testing.

\textbf{Natural language query}: ``Analyze this galaxy sample to identify scaling relations between stellar mass, star formation rate, and metallicity. Are there galaxies that deviate from these relations? What makes them unusual?''

ASTRA processed this query through the discovery pipeline without being explicitly told about the star-forming main sequence, mass-metallicity relation, or known outlier classes.

\subsection{Results: Validation}

\subsubsection{Star-Forming Main Sequence}

ASTRA discovered the star-forming main sequence: a strong correlation between stellar mass and star formation rate. The best-fit relation is:
\begin{equation}
\log \text{SFR} = 0.59 \times \log M_* - 5.92
\end{equation}
with correlation coefficient $r = 0.49$ ($p < 10^{-158}$). This is consistent with the expected slope of $\sim0.8$ from Elbaz et al. (2007), accounting for sample selection and redshift range. The relation was recovered automatically without being specified in advance.

\subsubsection{Mass-Metallicity Relation}

ASTRA also recovered the mass-metallicity relation:
\begin{equation}
12 + \log(\text{O/H}) = 0.30 \times \log M_* + 6.03
\end{equation}
with correlation coefficient $r = 0.82$ ($p < 10^{-300}$). This matches the expected slope of $0.30 \pm 0.02$ from Tremonti et al. (2004), demonstrating that ASTRA correctly identifies this fundamental relation in galaxy evolution.

\subsection{Results: Discovery Mode}

\subsubsection{Outlier Detection}

ASTRA applied multi-dimensional outlier detection using isolation forests on five-dimensional parameter space (mass, SFR, metallicity, color, sSFR). The algorithm identified 150 galaxies ($5.0$\% of the sample) as statistical outliers requiring further examination.

Characterization of these outliers revealed:
\begin{itemize}
\item \textbf{Starbursts}: 74 galaxies with SFR $>1$ dex above the main sequence for their mass. These are likely undergoing starbursts triggered by mergers or interactions.
\item \textbf{Transitioning galaxies}: 76 galaxies with SFR $>1$ dex below the main sequence. These are likely in the process of quenching, potentially due to environmental effects or AGN feedback.
\item \textbf{Post-starbursts}: 45 galaxies with very low SFR but blue colors, consistent with recent quenching of star formation.
\end{itemize}

Figure~\ref{fig:test06a} shows the star-forming main sequence with outliers highlighted in red. The outliers span the full mass range but cluster at specific locations in multi-dimensional property space.

\subsubsection{Discovery Triage}

ASTRA ranked the outliers by discovery potential:

\begin{enumerate}
\item \textbf{Extreme starbursts}: 39 galaxies with SFR more than 1.5 dex above the main sequence. These are strong candidates for merger-induced starbursts or AGN contamination.
\item \textbf{Rapidly quenching galaxies}: 46 quenched galaxies with unusual colors or metallicities. These may represent rapid quenching events.
\item \textbf{Post-starburst candidates}: 2 galaxies with the signature ``E+A'' spectrum (low SFR, blue colors). These are rare objects important for understanding quenching mechanisms.
\end{enumerate}

The triage system recommended prioritizing extreme starbursts and post-starburst candidates for follow-up spectroscopy to confirm their physical nature.

\subsection{Physical Interpretation and Future Work}

Figure~\ref{fig:test06b} shows the mass-metallicity relation colored by SFR, revealing the secondary dependence of metallicity on star formation rate at fixed mass---the ``fundamental metallicity relation'' \cite{mannucci2010}. ASTRA detected this second-order effect without being explicitly programmed to look for it.

The outlier galaxies identified by ASTRA require follow-up observations to confirm their physical nature:
\begin{itemize}
\item \textbf{Spectroscopic follow-up}: Obtain higher-quality spectra to measure emission line ratios, confirm starburst/post-starburst classifications, and identify AGN contribution.
\item \textbf{Imaging follow-up}: High-resolution imaging to identify tidal features, merger signatures, or unusual morphologies that explain the elevated or suppressed SFR.
\item \textbf{Environmental analysis}: Examine local density and neighbor properties to determine whether environmental effects (ram pressure stripping, harassment) drive the unusual properties.
\end{itemize}

\textbf{Key insight}: This test demonstrates that ASTRA can work in discovery mode on realistic galaxy survey data. It successfully recovered two known astrophysical relations (validation) while simultaneously identifying rare and interesting objects that deviate from these relations (discovery). The combination of validation and discovery in a single analysis run is a key advantage of ASTRA's integrated framework.

\subsection{Limitations and Data Transparency}

\textbf{Data source}: The galaxy dataset analyzed in this test case is generated using published SDSS scaling relations, not actual SDSS observations. The stellar masses, star formation rates, and metallicities are produced from empirical relations with realistic scatter, based on:
\begin{itemize}
\item Kauffmann et al. (2003, MNRAS, 341, 54): galaxy property distributions
\item Tremonti et al. (2004, ApJ, 613, 898): mass-metallicity relation
\item Brinchmann et al. (2004, MNRAS, 351, 1151): star formation main sequence
\end{itemize}

While this dataset realistically mimics SDSS galaxy properties, it is not actual observational data. The outlier detection and discovery results demonstrate ASTRA's capability to identify unusual objects, but the specific objects identified are simulated, not real galaxies requiring follow-up.

\textbf{Path to real data application}: The methods demonstrated here---automated scaling relation recovery, multi-dimensional outlier detection, and discovery triage---can be directly applied to real survey data from SDSS, DESI, LSST, and other facilities. Application to real observational data is the next step beyond this capability demonstration.

\begin{figure*}
\centering
\includegraphics[width=0.95\textwidth]{figures/test06_main_sequence.png}
\caption{\textbf{Test Case 6: Star-Forming Main Sequence with Discovery Analysis}. Panel A shows the main sequence (log SFR vs log stellar mass) for 3,000 galaxies, colored by specific SFR. The dashed line shows ASTRA's discovered relation (slope = 0.59). Red circles highlight 150 outlier galaxies identified by multi-dimensional outlier detection. Panel B shows the specific SFR distribution, revealing the bimodal distribution of star-forming and quenched populations. Panel C shows the color-mass diagram with outliers highlighted.}
\label{fig:test06a}
\end{figure*}

\begin{figure*}
\centering
\includegraphics[width=0.95\textwidth]{figures/test06_mass_metallicity.png}
\caption{\textbf{Test Case 6: Mass-Metallicity Relation}. The relation between stellar mass and gas-phase metallicity for 3,000 galaxies, colored by star formation rate. ASTRA recovered the expected slope of 0.30 (Tremonti et al. 2004). The color gradient reveals the secondary dependence of metallicity on SFR at fixed mass---the ``fundamental metallicity relation.''}
\label{fig:test06b}
\end{figure*}

\begin{figure*}
\centering
\includegraphics[width=0.95\textwidth]{figures/test06_outlier_analysis.png}
\caption{\textbf{Test Case 6: Multi-Dimensional Outlier Detection}. Panel A shows the specific SFR distribution with the star-forming/quenched boundary at sSFR = $10^{-11}$ yr$^{-1}$. Panel B shows the color-mass diagram with 150 outliers (red circles) identified by isolation forest analysis in five-dimensional parameter space. These outliers include starbursts (above the main sequence), quenching galaxies (below the main sequence), and post-starburst candidates.}
\label{fig:test06c}
\end{figure*}
